% This document is meant to be completed before watching the video:
%
% www.3blue1brown.com/?v=eulers-number
%
% The last row of the table can then be filled in a=e, 
% and the last column with ln(a)

\documentclass[9pt]{exam} 
\footer{}{}{}
\addpoints
\pointsinmargin

\usepackage{amsmath, amssymb}
\usepackage[a4paper, margin=20mm]{geometry} 
\usepackage{tcolorbox} 
\tcbset{colback=black!5!white}

\begin{document}

\textbf{Name:} \dotfill

\begin{questions}

%---------------------------------------------------------
\question[5] Develop the following expressions as linear combinations of $\ln(a)$, $\ln(b)$ and $\ln(c)$.
\begin{parts}
\vspace{4mm}
\part $\ln(ab) =$ \dotfill
\vspace{2mm}
\part $\ln\!\left(\dfrac{a^2 b^3}{c}\right) =$ \dotfill
\part $\ln\!\left(\dfrac{a\sqrt{b}}{c^3}\right) =$ \dotfill
\end{parts}

%---------------------------------------------------------
\question[5] Compute the following derivatives.
\begin{parts}
\part $\dfrac{d}{dx}(\ln x) =$ \dotfill
\part $\dfrac{d}{dx}(x \ln x) =$ \dotfill
\part $\dfrac{d}{dx}(\ln(3x^2 + 1)) =$ \dotfill
\fillwithdottedlines{14mm}
\end{parts}


%---------------------------------------------------------
\question[2] Reminder: $(e^x)' = e^x$. Use this to compute the derivative of $e^{cx}$, where $c \in \mathbb{R}$ is a constant.
\fillwithdottedlines{21mm}


%---------------------------------------------------------
\question[4] Which of the following functions models the size of a population which \textbf{doubles} every day after $x$ days?

\begin{oneparchoices}
\choice $f(x) = 2x$
\choice $f(x) = x^2$
\choice $f(x) = 2^x$
\choice $f(x) = x + 2$
\end{oneparchoices}

\begin{minipage}[t]{0.5\textwidth}
Complete the table for your chosen function.

\vspace{4mm}

{\Large
\renewcommand{\arraystretch}{1.2}
\begin{tabular}{c|cccccc}
$x$ & 0 & 1 & 2 & 3 & 4 & 10 \\ \hline
$f(x)$ & & & & & & 
\end{tabular}
}

\end{minipage}
\hfill
\begin{minipage}[t]{0.5\textwidth}

Solve algebraically for $f(x) = 900$. \hfill \emph{CRM p.14}
\fillwithdottedlines{21mm}

\end{minipage}

%---------------------------------------------------------
\question[2] Complete the following table for the expression $\dfrac{a^{dt} - 1}{dt}$,
using $dt = 1,\; 0.01,\; 0.00001$ and $a = 2, 3, 8$.

\begin{center} \Large
$\displaystyle \dfrac{a^{dt} - 1}{dt} \qquad $
\renewcommand{\arraystretch}{1.2}
\begin{tabular}{c|cccc|c}
$a$ & $dt=1$ & $0.01$ & $0.00001$ & & \hspace{30mm} \\ \hline
2 &  &  &  &  & \\ \hline
3 &  &  &  &  & \\ \hline
8 &  &  &  &  & \\ \hline
 &  &  &  &  &
\end{tabular}
\end{center}


\vfill

\question[0] \textbf{Bonus.} For $f(x) = x \ln (x)$, the equation of the tangent line at $x = e$ is \dotfill \emph{\footnotesize(work on the back).}

\end{questions}

%---------------------------------------------------------

\begin{tcolorbox}
\textbf{Name of marker:}

\begin{center}
\gradetable[h][questions]
\end{center}

\end{tcolorbox}

\end{document}

\begin{center}
\renewcommand{\arraystretch}{1.2}
\begin{tabular}{c|cccc|c}
$a$ & $\Delta x=1$ & $0.01$ & $0.0001$ & $0.000001$ & \hspace{30mm} \\ \hline
2 & 1.0000 & 0.693 & 0.6931 & 0.69315 & \\
3 & 2.0000 & 1.099 & 1.0986 & 1.09861 & \\
8 & 7.0000 & 2.079 & 2.0794 & 2.07944 &
\end{tabular}
\end{center}